% ── Tự động tạo bởi scripts/analyze_rtt_latency.py ──────────────────────
\begin{figure}[H]
\centering
\begin{tikzpicture}
\begin{axis}[
    ybar,
    width=0.88\linewidth,
    height=6cm,
    bar width=12pt,
    xlabel={Điều kiện tải mạng},
    ylabel={RTT (ms)},
    symbolic x coords={Bình thường, Tải cao 100 Mbit/s, Tải cao 150 Mbit/s, Tải cao 200 Mbit/s},
    xtick=data,
    xticklabel style={font=\small},
    ymin=0, ymax=110,
    ytick={0,20,40,60,80,100},
    ymajorgrids=true,
    major grid style={line width=0.4pt, draw=gray!40},
    legend style={at={(0.98,0.98)}, anchor=north east, font=\small,
                  draw=gray!40, fill=white},
    legend cell align=left,
    axis lines=left,
    enlarge x limits=0.20,
]
% Trung bình (TB)
\addplot[fill=blue!60, draw=blue!80] coordinates {
    (Bình thường, 22.6)
    (Tải cao 100 Mbit/s, 27.4)
    (Tải cao 150 Mbit/s, 38.7)
    (Tải cao 200 Mbit/s, 50.2)
};
\addlegendentry{TB}
% P50 (Trung vị)
\addplot[fill=green!60, draw=green!80!black] coordinates {
    (Bình thường, 19.0)
    (Tải cao 100 Mbit/s, 23.5)
    (Tải cao 150 Mbit/s, 33.3)
    (Tải cao 200 Mbit/s, 52.8)
};
\addlegendentry{P50}
% P95
\addplot[fill=orange!70, draw=orange!90] coordinates {
    (Bình thường, 52.4)
    (Tải cao 100 Mbit/s, 58.3)
    (Tải cao 150 Mbit/s, 82.3)
    (Tải cao 200 Mbit/s, 93.4)
};
\addlegendentry{P95}
% P99
\addplot[fill=red!60, draw=red!80] coordinates {
    (Bình thường, 80.0)
    (Tải cao 100 Mbit/s, 88.0)
    (Tải cao 150 Mbit/s, 95.8)
    (Tải cao 200 Mbit/s, 98.1)
};
\addlegendentry{P99}
\end{axis}
\end{tikzpicture}
\caption{RTT (TB / P50 / P95 / P99) theo điều kiện tải mạng Wi-Fi giữa ứng dụng VR (Meta Quest~3) và Ros2 (PC)}
\label{fig:rtt-bar-chart}
\end{figure}

\begin{table}[H]
\centering
\caption{Thống kê RTT giữa ứng dụng VR và RosBridge WebSocket theo điều kiện tải mạng}
\label{tab:rtt-latency}
\small
\begin{tblr}{
    colspec = {|X[3.2,l]|X[1,c]|X[1.2,c]|X[1.2,c]|X[1.2,c]|X[1.2,c]|X[1.2,c]|X[1.2,c]|},
    hlines,
    row{1} = {font=\bfseries, bg=gray!20},
}
Điều kiện & N & Min (ms) & TB (ms) & Trung vị (ms) & P95 (ms) & P99 (ms) & Max (ms) \\
Bình thường & 1426 & 0{,}8 & 22{,}6 & 19{,}0 & 52{,}4 & 80{,}0 & 97{,}3 \\
Tải cao 100 Mbit/s & 1001 & 0{,}7 & 27{,}4 & 23{,}5 & 58{,}3 & 88{,}0 & 99{,}4 \\
Tải cao 150 Mbit/s & 813 & 0{,}4 & 38{,}7 & 33{,}3 & 82{,}3 & 95{,}8 & 99{,}4 \\
Tải cao 200 Mbit/s & 955 & 0{,}0 & 50{,}2 & 52{,}8 & 93{,}4 & 98{,}1 & 99{,}8 \\
\end{tblr}
\end{table}

Kết quả đo đạc RTT được trình bày tại Bảng~\ref{tab:rtt-latency} và Hình~\ref{fig:rtt-bar-chart}.
Đây là thời gian khứ hồi (Round-Trip Time) tính từ khi Ros2 publish topic
\texttt{/benchmark} cho đến khi nhận được gói phản hồi \texttt{/ack}
từ ứng dụng VR thông qua kết nối WebSocket RosBridge tại cổng~9090.

\textbf{Điều kiện mạng bình thường.}
Trong điều kiện không có tải mạng tải cao, RTT trung bình đạt
\textbf{22{,}6\,ms} (trung vị 19{,}0\,ms),
dao động từ 0{,}8\,ms đến 97{,}3\,ms với độ lệch chuẩn
13{,}9\,ms.
Ngưỡng P95 ở mức \textbf{52{,}4\,ms} cho thấy 95\% các gói tin
đến trong vòng 52\,ms -- thỏa mãn tốt yêu cầu điều khiển
theo thời gian thực với chu kỳ 100\,ms của Ros2-Control.

\textbf{Ảnh hưởng của tải mạng tải cao.}
Khi áp tải mạng 100\,Mbit/s, RTT trung bình tăng lên
\textbf{27{,}4\,ms} (+21{,}4\%
so với điều kiện bình thường), P99 đạt 88{,}0\,ms.
Ở mức 150\,Mbit/s, RTT trung bình là \textbf{38{,}7\,ms}
(+71{,}5\%),
và tại 200\,Mbit/s lên tới \textbf{50{,}2\,ms}
(+122{,}3\%).
Mặc dù giá trị trung bình tăng đáng kể khi tải mạng vượt 150\,Mbit/s,
P99 vẫn nằm dưới 100\,ms ở tất cả các điều kiện kiểm thử,
đồng nghĩa với việc hệ ít hơn 1\% số gói tin bị trễ vượt mức cho phép
ngay cả trong điều kiện tải mạng cực cao.

\textbf{Nhận xét chung.}
RTT giữa ứng dụng VR và Ros2 đạt mức thấp nhờ kết nối mạng LAN nội bộ
giữa kính Meta Quest~3 và máy PC chạy Ros2 (cùng mạng Wi-Fi 5\,GHz).
Độ trễ \textasciitilde23\,ms ở điều kiện bình thường
nhỏ hơn đáng kể so với chu kỳ điều khiển 100\,ms của Ros2-Control,
đảm bảo rằng lệnh từ ứng dụng VR được truyền đến robot trong cùng chu kỳ điều khiển.
